\section{Lecture 13: Generalities on Analytic Rings --- Completion, Colimits, and Frobenius}
\label{lec:13}

Having spent the previous lectures building the solid theory and its Huber-pair
variants, and having seen Scholze introduce the liquid and gaseous structures
that also work archimedeanly (\cref{lec:12}), Clausen now steps back to develop
some general operations on the \emph{category} of analytic rings. The immediate
motivation is that manipulating the new (liquid, gaseous, $\dots$) examples is
much easier with a good handle on how analytic rings are constructed, mapped,
and glued; but the material should be recorded in any case. There are three
parts.

\begin{itemize}
\item A \emph{completion} functor turning a ``pre-analytic ring'' (the axioms of
an analytic ring minus the requirement that the unit be complete) into an
honest analytic ring, and its derived ($\mathbb{E}_\infty$/animated)
enhancement.
\item \emph{Colimits} of analytic rings: computed by taking the colimit of
pre-analytic rings and completing. Filtered colimits need no completion;
pushouts (completed tensor products) are the substantive case.
\item A theorem that \emph{Frobenius is a map of analytic rings}, via a
free-module criterion for maps of analytic rings and a Tate-cohomology
computation that is special to the light setting. As a payoff one gets derived
symmetric powers on the derived category, hence a good theory of animated
analytic rings.
\end{itemize}

The lecture closes with a general ``killing an algebra object'' construction
producing left adjoints to reflective subcategories, which uniformly recovers
solidification, localization, and completion.

Throughout, ``ring'' means commutative ring, everything lives in light condensed
abelian groups, and we use freely the abstract notion of analytic ring and its
derived category (\cref{lec:8}), the solid theory (\cref{lec:5,lec:6}), and its
localization over the (valuative) spectrum (\cref{lec:9,lec:10}). We write $\ur R$
for the underlying condensed ring of an analytic ring $R$ (the ``triangle''
notation of \cref{lec:1,lec:8}).

\subsection{Analytic and pre-analytic rings}

We first recall the definition in the form Clausen will use, then relax it.

\begin{definition}[Analytic ring]
\label{def:13:analytic-ring}
An \emph{analytic ring} is a pair $(\ur R, \Mod_R)$ where $\ur R$ is a condensed
(commutative) ring and
\begin{equation}
\label{eq:13:mod-subcat}
\Mod_R\ \subseteq\ \Mod_{\ur R}
\end{equation}
is a full subcategory of condensed $\ur R$-modules, closed under limits,
colimits and extensions, such that
\begin{equation}
\label{eq:13:ext-axiom}
\iExt^{i}_{\ur R}(M,N)\ \in\ \Mod_R\qquad\text{for all }N\in\Mod_R,\ M\in\Mod_{\ur R},\ i\ge 0,
\end{equation}
and such that, finally,
\begin{equation}
\label{eq:13:star}
\ur R\ \in\ \Mod_R.\tag{$\star$}
\end{equation}
We call \eqref{eq:13:star} the \emph{completeness} axiom.
\end{definition}

\begin{definition}[Pre-analytic ring]
\label{def:13:pre-analytic-ring}
A \emph{pre-analytic ring} is a pair $(\ur R, \Mod_R)$ as in
\cref{def:13:analytic-ring} except that we drop \eqref{eq:13:star}: the unit
object $\ur R$ need not be complete.
\end{definition}

\begin{remark}[Why relax the definition]
\label{rmk:13:why-pre-analytic}
The completeness axiom \eqref{eq:13:star} is exactly the clause that is awkward
to verify directly on an interesting condensed ring $\ur R$: it forces one to
control $\Ext$-groups in $\Mod_{\ur R}$, i.e.\ to understand the homological
algebra of condensed $\ur R$-modules before any computation can begin --- hard,
for instance, for a liquid base like a ring of convergent Laurent series. In
practice one instead maps a \emph{simpler} ring into $\ur R$ (e.g.\
$\mathbb{Z}[T,T^{-1}]$ for the Laurent case) on which a pre-analytic structure is
easy to write down, and then \emph{completes}. We saw the same phenomenon for
adic spaces: the module category attached to a rational open was first specified
by imposing conditions --- a pre-analytic ring --- and one then had to change
$\ur R$ to (essentially) the value of Huber's structure sheaf to make
\eqref{eq:13:star} hold. The completion procedure below is what carries out this
last step.
\end{remark}

\begin{definition}[Maps]
\label{def:13:maps}
A \emph{map of pre-analytic rings} $(\ur R,\Mod_R)\to(\ur S,\Mod_S)$ is a map of
underlying condensed rings $\ur R\to\ur S$ such that restriction of scalars sends
$\Mod_S$ into $\Mod_R$: for $N\in\Mod_S$, the underlying $\ur R$-module lies in
$\Mod_R$. Equivalently, writing $\widehat{(-)}_R\colon\Mod_{\ur R}\to\Mod_R$ and
$\widehat{(-)}_S\colon\Mod_{\ur S}\to\Mod_S$ for the completion (localization)
functors, the condition says that base change followed by $S$-completion
descends to a functor
\begin{equation}
\label{eq:13:base-change-map}
-\otimes_R S\ \colon\ \Mod_R\ \longrightarrow\ \Mod_S,
\end{equation}
which is then automatically symmetric monoidal. (Concretely, this uses the
universal property of the localization $\widehat{(-)}_R$: a map that becomes an
isomorphism after $R$-completion also becomes one after $S$-completion.)
\end{definition}

Since analytic rings are pre-analytic rings, there is a fully faithful inclusion
\begin{equation}
\label{eq:13:ff-inclusion}
\{\text{analytic rings}\}\ \hookrightarrow\ \{\text{pre-analytic rings}\}.
\end{equation}

\subsection{Completion}

\begin{proposition}[Completion is a left adjoint]
\label{prop:13:completion}
The inclusion \eqref{eq:13:ff-inclusion} has a left adjoint. On objects it sends
a pre-analytic ring $(\ur R, \Mod_R)$ to
\begin{equation}
\label{eq:13:completion-formula}
(\ur R, \Mod_R)\ \longmapsto\ (\widehat{\ur R},\, \Mod_R),
\end{equation}
where $\widehat{\ur R}$ is the completion of the unit --- the image of $\ur R$
under the left adjoint $\widehat{(-)}_R\colon\Mod_{\ur R}\to\Mod_R$ to the
inclusion --- and $\Mod_R$ is the same module category, now viewed inside
$\Mod_{\widehat{\ur R}}$.
\end{proposition}

\begin{proof}[Proof sketch]
There are three things to check: that $\widehat{\ur R}$ is a condensed ring;
that $\Mod_R$ is a full subcategory of $\Mod_{\widehat{\ur R}}$; and that
$(\widehat{\ur R},\Mod_R)$ satisfies the axioms of a (pre-)analytic ring,
including \eqref{eq:13:star}.

\emph{$\widehat{\ur R}$ is a commutative condensed ring.} This is completely
formal. The completion functor
$\widehat{(-)}_R\colon\Mod_{\ur R}\to\Mod_R$ is symmetric monoidal and
colimit-preserving, and its right adjoint is the inclusion
$\Mod_R\hookrightarrow\Mod_{\ur R}$. The right adjoint of a symmetric monoidal
functor is lax symmetric monoidal, hence carries the unit object of $\Mod_R$ to a
commutative algebra object of $\Mod_{\ur R}$; that algebra is precisely
$\widehat{\ur R}$. A fortiori $\widehat{\ur R}$ is a condensed commutative ring.
By the same formal fact, the right adjoint sends any object of $\Mod_R$ to a
module over $\widehat{\ur R}$, so the inclusion factors
\begin{equation}
\label{eq:13:factoring}
\Mod_R\ \hookrightarrow\ \Mod_{\widehat{\ur R}}\ \longrightarrow\ \Mod_{\ur R},
\end{equation}
the second map being restriction of scalars along $\ur R\to\widehat{\ur R}$.
% board 8 @ 00:14:50 --- the factoring of Mod_R through modules over the completed unit
\[
\begin{tikzcd}
\Mod_R \arrow[rr, hook] \arrow[dr] & & \Mod_{\ur R} \\
& \Mod_{\widehat{\ur R}} \arrow[ur] &
\end{tikzcd}
\]

\emph{$\Mod_R\hookrightarrow\Mod_{\widehat{\ur R}}$ is fully faithful.} For
$M,N\in\Mod_R$ one compares $\Homop_{\ur R}(M,N)$ with
$\Homop_{\widehat{\ur R}}(M,N)$. Since $N$ is an $\widehat{\ur R}$-module,
\[
\Homop_{\widehat{\ur R}}(M,N)=\Homop_{\widehat{\ur R}}\big(\widehat{\ur R}\otimes_{\ur R}M,\,N\big);
\]
as $N$ is complete one may apply $R$-completion inside, and, completion being
symmetric monoidal and idempotent, this collapses to $\Homop_{\widehat{\ur R}}(M,N)=\Homop_{\ur R}(M,N)$. Hence
$\Mod_R\hookrightarrow\Mod_{\widehat{\ur R}}$ is full and faithful.

\emph{The axioms.} Closure under limits and colimits is automatic since the
forgetful functor $\Mod_{\widehat{\ur R}}\to\Mod_{\ur R}$ preserves them. The
completeness axiom \eqref{eq:13:star} holds by construction:
$\widehat{\ur R}\in\Mod_R$. Closure under extensions and the internal-$\Ext$
axiom \eqref{eq:13:ext-axiom} are again $\Ext$-computations; the fully-faithful
argument upgrades to them, but only \emph{provided the derived completion of
$\ur R$ agrees with the underived completion} $\widehat{\ur R}$ --- see
\cref{rmk:13:derived-needed}.
\end{proof}

\begin{remark}[Why one is forced into the derived setting]
\label{rmk:13:derived-needed}
The $\Ext$-part of \cref{prop:13:completion} is subtle. The derived completion
$\widehat{\ur R}$ need not sit in degree $0$: its $H_0$ is the underived
completion, but it can carry higher homology (this is exactly the phenomenon
behind the non-sheafy adic spaces of \cref{lec:10}). The formal argument for the
$\Ext$-axioms goes through cleanly \emph{if} one replaces $\widehat{\ur R}$ by
its derived version and uses a notion of \emph{derived analytic ring}. So it is
convenient to introduce that notion in the middle of this proof.
\end{remark}

\subsection{Derived analytic rings and the abelian comparison}

\begin{definition}[Analytic $\mathbb{E}_\infty$-ring]
\label{def:13:derived-analytic-ring}
An \emph{analytic $\mathbb{E}_\infty$-ring} (or \emph{derived analytic ring}) is a
pair $(\ur R, D_{\ge 0}(R))$ where $\ur R$ is a condensed connective
$\mathbb{E}_\infty$-ring and
\begin{equation}
\label{eq:13:derived-subcat}
D_{\ge 0}(R)\ \subseteq\ D_{\ge 0}(\ur R)
\end{equation}
is a full subcategory (of the connective part of the derived category), closed
under all limits and colimits, such that
\begin{equation}
\label{eq:13:derived-hom-axiom}
\iHom_{\ur R}(M,N)\ \in\ D_{\ge 0}(R)\qquad\text{for }N\in D_{\ge 0}(R),\ M\in D_{\ge 0}(\ur R).
\end{equation}
It is a \emph{pre-analytic} $\mathbb{E}_\infty$-ring if we do not require the
unit; an (honest) \emph{analytic} one if in addition $\ur R\in D_{\ge 0}(R)$.
Only the connective part is specified; note that closure under the loop functor
in $D_{\ge 0}$ --- shift down and truncate --- forces the homology-wise
description below.
\end{definition}

\begin{proposition}[Derived vs.\ abelian structures]
\label{prop:13:derived-abelian-bijection}
For any connective condensed $\mathbb{E}_\infty$-ring $\ur R$ there is a bijection
\begin{equation}
\label{eq:13:bijection}
\begin{aligned}
&\{\text{pre-analytic ring structures on }\ur R\}\\
&\hspace{4em}\xrightarrow{\ \sim\ }
\{\text{pre-analytic ring structures on }\pi_0\ur R\ (\text{old, abelian sense})\}.
\end{aligned}
\end{equation}
In the forward direction one restricts to the heart,
\begin{equation}
\label{eq:13:forward}
D_{\ge 0}(R)\ \longmapsto\ D_{\ge 0}(R)\cap D_{=0}(\ur R)\ =\ D_{=0}(\pi_0\ur R),
\end{equation}
and in the backward direction one takes homology termwise,
\begin{equation}
\label{eq:13:backward}
\Mod_R\ \longmapsto\ \{M\in D_{\ge 0}(\ur R)\ :\ H_i M\in\Mod_R\ \forall\, i\ge 0\}.
\end{equation}
An analytic $\mathbb{E}_\infty$-ring structure maps to an analytic (abelian) one;
the converse holds when $\ur R=\pi_0\ur R$, and in general the analytic
structures on $\ur R$ correspond to those abelian structures on $\pi_0\ur R$ for
which every homotopy group $\pi_i\ur R$ (not only $\pi_0$) is complete.
\end{proposition}

\begin{remark}[Consistency in the classical case]
\label{rmk:13:consistency}
When $\ur R=\pi_0\ur R$ is concentrated in degree $0$ (a classical ring), the two
notions overlap and \cref{prop:13:derived-abelian-bijection} says they agree: the
derived definition of a (pre-)analytic ring structure matches the naive
abelian-level definition of \cref{lec:8}.
\end{remark}

\Cref{prop:13:derived-abelian-bijection} finishes the proof of
\cref{prop:13:completion}. There is no obstruction to proving the completion
proposition in the derived setting (\cref{rmk:13:derived-needed}); one then
transports the resulting analytic-ring structure on the derived $\widehat{\ur R}$
down to its $\pi_0$ by the recipe \eqref{eq:13:forward}. The proof of
\cref{prop:13:derived-abelian-bijection} itself was given in an earlier iteration
of the course \cite{clausen2020analytic}.

\begin{remark}[Warning: the derived category can change]
\label{rmk:13:derived-cat-warning}
An analytic ring structure on $\pi_0\ur R$ produces both an abelian category and,
via \eqref{eq:13:backward}, a derived category. If
$(\widehat{\ur R},\Mod_R)$ is the (abelian) completion of $(\ur R,\Mod_R)$ from
\cref{prop:13:completion}, it is \emph{not} in general true that
\begin{equation}
\label{eq:13:derived-cat-mismatch}
D(\widehat{\ur R},\Mod_R)\ =\ D(\ur R,\Mod_R);
\end{equation}
this holds only if one uses \emph{derived} completion on the left-hand side and
that derived completion happens to live in degree $0$. When the derived
completion has higher homology one should genuinely work with a derived analytic
ring rather than its $\pi_0$. This is the mechanism behind the derived structure
sheaves of the exotic rational localizations of \cref{lec:10}: even though one has
an underlying ordinary analytic ring by taking $\pi_0$, the correct object is the
derived analytic ring. Finally, maps of derived analytic rings are the same as
maps of the underlying derived rings, which on hearts recover maps in the old
abelian sense.
\end{remark}

\subsection{Colimits of analytic rings}

The completion functor of \cref{prop:13:completion} lets us compute colimits.

\begin{proposition}[Colimits: colimit of pre-analytic rings, then complete]
\label{prop:13:colimits}
Colimits in analytic rings are computed by taking the colimit in pre-analytic
rings and completing. The colimit of a diagram $(\ur R_i,\Mod_{R_i})_{i\in I}$ of
pre-analytic rings is
\begin{equation}
\label{eq:13:colimit-formula}
\varinjlim_{i\in I}(\ur R_i,\Mod_{R_i})\ =\ \Big(\varinjlim_i \ur R_i,\ \Mod_{\varinjlim R_i}\Big),
\end{equation}
where $\varinjlim_i\ur R_i$ is the colimit in condensed commutative rings and
\begin{equation}
\label{eq:13:colimit-modules}
\Mod_{\varinjlim R_i}\ =\ \big\{M\in\Mod_{\varinjlim \ur R_i}\ :\ \forall i,\ M|_{\ur R_i}\in\Mod_{R_i}\big\},
\end{equation}
i.e.\ the intersection over $i$ of the conditions of lying in $\Mod_{R_i}$ after
restriction of scalars.
\end{proposition}

That \eqref{eq:13:colimit-formula} is the colimit in pre-analytic rings is
immediate from the definition of a map (\cref{def:13:maps}): a map out of the
colimit is a map of the underlying rings out of $\varinjlim_i\ur R_i$ with the
correct completeness property, and \eqref{eq:13:colimit-modules} is tailored to
encode exactly that. Checking the pre-analytic axioms is elementary (limits and
colimits because the relevant forgetful functors preserve them; the $\Ext$-axiom
directly). That colimits of analytic rings are then computed by completing is
formal, the completion being a left adjoint.

\begin{example}[Filtered colimits need no completion]
\label{ex:13:filtered}
For a \emph{filtered} diagram the completion is unnecessary:
$(\varinjlim_i\ur R_i,\Mod_{\varinjlim R_i})$ is already complete. Indeed, to
test \eqref{eq:13:star} one must check that the unit $\varinjlim_i\ur R_i$ lies in
$\Mod_{R_i}$ for every $i$; but from $i$ onwards it is a filtered colimit of
objects of $\Mod_{R_i}$ (cofinally), and $\Mod_{R_i}$ is closed under filtered
colimits, so it lies there. Hence the pre-analytic colimit is already an analytic
ring. For instance, for a discrete commutative ring $R$ written as a filtered
colimit $R=\varinjlim_i R_i$,
\begin{equation}
\label{eq:13:solid-filtered}
\Solid_R\ =\ \varinjlim_{i\in I}\Solid_{R_i}\qquad(\text{colimit in analytic rings}),
\end{equation}
as we already saw in \cref{lec:11}.
\end{example}

\begin{remark}[Sifted suffices; the only difficulty is finite]
\label{rmk:13:sifted}
One need not say ``filtered'' in \cref{ex:13:filtered}: \emph{sifted} colimits are
just as easy. Since general colimits decompose into sifted colimits and finite
coproducts (just as they decompose into filtered colimits and pushouts), the
completion functor for a relative tensor product is the same as for the absolute
one, and the only genuine difficulty is the finite case. In other words
$A\otimes_R B$ is no harder than $A\otimes B$; the substance is entirely in the
pushout.
\end{remark}

\begin{example}[Pushouts and completed tensor products]
\label{ex:13:pushouts}
For a span $A\leftarrow R\rightarrow B$ of analytic rings, the pushout is the
completion of the pre-analytic ring with underlying ring $\ur A\otimes_{\ur R}\ur B$
and module category
\begin{equation}
\label{eq:13:pushout-modules}
\big\{\,M\in\Mod_{\ur A\otimes_{\ur R}\ur B}\ :\ M|_{\ur A}\in\Mod_A\ \text{and}\ M|_{\ur B}\in\Mod_B\,\big\}.
\end{equation}
Here the completion is essential, since $\ur A\otimes_{\ur R}\ur B$ has no reason
to lie in $\Mod_A$ or $\Mod_B$.
\end{example}

\begin{remark}[The pushout completion iterates]
\label{rmk:13:pushout-iteration}
The completion functor in \cref{ex:13:pushouts} a priori involves \emph{iterating}
$A$-completion (of the underlying $\ur A$-module) and $B$-completion (of the
underlying $\ur B$-module), alternately, along a countable sequential colimit. In
practice it is usually not so bad, but that is what one must do in general. The
same discussion applies verbatim in the derived/animated context, interpreting
$\otimes$ as $\otimes^{\mathbb L}$ and $\Mod$ as $D_{\ge 0}$.
\end{remark}

\begin{example}[Rational opens as pushouts]
\label{ex:13:rational-opens}
In the solid $(R,R^+)$-theory of \cref{lec:9}, a rational open $U\subseteq\Spv(R,R^+)$
gives a new analytic ring $(\mathcal O(U),\mathcal O^+(U))$: the completion of the
pre-analytic structure on $R$ in which one enforces the relations describing $U$
--- for $U=U(f_1,\dots,f_n/g)$, that $g$ act invertibly and each $f_i/g$ be a
solid variable on all modules. Pushouts then compute intersections:
\begin{equation}
\label{eq:13:pushout-intersection}
(\mathcal O(U),\mathcal O^+(U))\ \otimes_{(R,R^+)}\ (\mathcal O(V),\mathcal O^+(V))\ =\ (\mathcal O(U\cap V),\mathcal O^+(U\cap V)),
\end{equation}
directly from \eqref{eq:13:pushout-modules} (in the discrete case
$\otimes_{(R,R^+)}$ is literally $\otimes_R$). Geometrically these pushouts are
the important construction: they are what should correspond to pullbacks (here,
intersections of opens), and the need to complete is the additional subtlety of
analytic geometry over ordinary algebraic geometry.
\end{example}

\subsection{Frobenius is a map of analytic rings}

Fix a prime $p$.

\begin{theorem}[Frobenius]
\label{thm:13:frobenius}
Let $R$ be an analytic ring. Then the Frobenius map
\begin{equation}
\label{eq:13:frobenius-map}
\Frob\colon\ \ur R\ \longrightarrow\ \ur R/p,\qquad x\longmapsto x^p,
\end{equation}
is a map of analytic rings $R\to R/p$. Here $\ur R/p$ means $\widehat{\ur R}/p$,
and $R/p$ carries the induced analytic ring structure
\begin{equation}
\label{eq:13:Rmodp}
R/p\ =\ \big(\ur R/p,\ \Mod_{\ur R/p}(\Mod_R)\big),
\end{equation}
i.e.\ a complete $\ur R/p$-module is an $\ur R/p$-module that is complete when
viewed as an $\ur R$-module.
\end{theorem}

To lighten notation one may assume $\ur R$ has characteristic $p$ (is an
$\mathbb{F}_p$-algebra), so that $\Frob\colon\ur R\to\ur R$. By
\eqref{eq:13:Rmodp} the content is:
\begin{equation}
\label{eq:13:frobenius-goal}
M\in\Mod_R\ \Longrightarrow\ \Frob_* M\in\Mod_R,
\end{equation}
where $\Frob_*M$ is $M$ with $\ur R$-action restricted along Frobenius. This is
not obvious: the axioms of an analytic ring are all categorical closure
properties in linear algebra (limits, colimits, $\Ext$), and say nothing about
Frobenius. One needs another handle on maps of analytic rings.

\subsubsection*{A free-module criterion for maps}

\begin{proposition}[Necessary data]
\label{prop:13:free-module-data}
If $(\ur R,\Mod_R)\to(\ur S,\Mod_S)$ is a map of analytic rings, then for every
light profinite set $T$ one gets an $\ur R$-linear map of free complete modules
\begin{equation}
\label{eq:13:free-module-map}
R[T]\ \longrightarrow\ S[T],
\end{equation}
functorial in $T$. (By hypothesis the free complete $S$-module $S[T]$ lies in
$\Mod_R$ after restriction; hence the map $\ur R[T]\to S[T]$ from the uncompleted
free module factors through the completion $R[T]$. It suffices to have this for a
generating class of $T$, e.g.\ the Cantor set.) Moreover the maps
\eqref{eq:13:free-module-map} are compatible with maps into the ground ring: for
$t\colon T\to\ur R$ a map of condensed sets, the square
\begin{equation}
\label{eq:13:compatibility-square}
\begin{tikzcd}[column sep=large]
R[T] \arrow[r] \arrow[d] & S[T] \arrow[d] \\
\ur R \arrow[r] & \ur S
\end{tikzcd}
\end{equation}
commutes, where $R[T]\to\ur R$ is the extension of $t$ (using completeness of
$\ur R$) and the right and lower maps are the evident ones.
\end{proposition}

\begin{lemma}[Sufficiency of the free-module data]
\label{lem:13:sufficiency}
Conversely, if $\ur R\to\ur S$ is a map of condensed rings for which there exist
$\ur R$-linear maps $R[T]\to S[T]$, functorial in $T$ and satisfying the
compatibility of \eqref{eq:13:compatibility-square}, then $\ur R\to\ur S$ is a
map of analytic rings $R\to S$.
\end{lemma}

The point of \cref{lem:13:sufficiency} is that one need not build the base-change
functor $\Mod_S\to$ ($S$-modules from $R$-modules) by hand: producing the maps on
free objects, plus the one compatibility, is enough.

\begin{proof}[Idea]
$\Mod_R$ is monadic over condensed sets: the forgetful functor
$\Mod_R\to\Mod_{\ur R}\to(\text{condensed sets})$ is the right adjoint of a
localization followed by a forgetful functor, so it satisfies the hypotheses of
the Barr--Beck monadicity theorem, and $\Mod_R$ is the category of algebras over
the free-complete-$R$-module monad. To show every $S$-module is an $R$-module it
suffices to produce a map of monads from the free-$R$-module monad to the
free-$S$-module monad; its first datum is exactly a map of free objects
\eqref{eq:13:free-module-map}, and unwinding the compatibility with the monad
structure reduces to the single commuting square
\eqref{eq:13:compatibility-square}. The full argument is in the course notes
\cite{clausen2020analytic}.
\end{proof}

By \cref{prop:13:free-module-data} and \cref{lem:13:sufficiency}, to prove
\cref{thm:13:frobenius} it suffices to construct, for every light profinite set
$T$, a \emph{Frobenius-semilinear} map $R[T]\to R[T]$, functorial in $T$ and
compatible as in \eqref{eq:13:compatibility-square}.

\subsubsection*{Recollection on Frobenius}

Let $C_p$ denote the cyclic group of order $p$, and write
$\widehat H^{0}(C_p,-)=\pi_0\big((-)^{tC_p}\big)$ for degree-zero Tate cohomology.

\begin{recollection}[The abelian Frobenius]
\label{rec:13:frobenius}
For any abelian group $M$ there is a natural map
\begin{equation}
\label{eq:13:tate-frobenius}
M\ \longrightarrow\ \widehat H^{0}\big(C_p,\,M^{\otimes p}\big)\ =\ \coker\!\big(N\colon (M^{\otimes p})_{C_p}\to (M^{\otimes p})^{C_p}\big)\ =\ \pi_0\big(M^{\otimes p}\big)^{tC_p},
\end{equation}
where $C_p$ permutes the tensor factors of $M^{\otimes p}$ and $N$ is the norm
(transfer) map. It is induced by the diagonal
\begin{equation}
\label{eq:13:diagonal}
m\ \longmapsto\ [\,m\otimes\cdots\otimes m\,]\ \in\ (M^{\otimes p})^{C_p}\big/N\big((M^{\otimes p})_{C_p}\big),
\end{equation}
which is $C_p$-fixed. Although $m\mapsto m^{\otimes p}$ is not additive, the
cross-terms in the image of $m+n$ are each norms, hence vanish in the Tate
quotient, so \eqref{eq:13:tate-frobenius} is a group \emph{homomorphism}. (A nice
exercise, if one has not done it.)
\end{recollection}

If $M=R$ is a commutative ring, composing with the ($C_p$-equivariant)
multiplication $R^{\otimes p}\to R$, on whose target $C_p$ acts trivially, gives
\begin{equation}
\label{eq:13:ring-frobenius}
R\ \longrightarrow\ \pi_0\big(R^{\otimes p}\big)^{tC_p}\ \longrightarrow\ \pi_0\, R^{tC_p}\ =\ R/p,\qquad m\longmapsto m^p,
\end{equation}
the usual Frobenius: as $C_p$ acts trivially on $R$, the norm is multiplication
by $p$ and $R^{tC_p}=R/p$. If $M$ is an $R$-module, the map $M\to\pi_0(M^{\otimes p})^{tC_p}$ is a
map of abelian groups, and in fact a map of $R$-modules once one Frobenius-twists
the right-hand side --- the abelian-group Frobenius is linear over the ring
Frobenius.

\subsubsection*{Reduction to a key lemma}

This suggests applying the abelian Frobenius to the free module $R[T]$, viewed
merely as a sheaf of abelian groups. For a light profinite set $T$ one gets a
sequence
% board 3 @ 01:11:22 --- applying the abelian Frobenius to R[T]; the diagonal Delta is an iso
\begin{equation}
\label{eq:13:frobenius-RT}
\begin{tikzcd}[column sep=large]
R[T] \arrow[r] & \pi_0\!\big(R[T]^{\otimes p}\big)^{tC_p} \arrow[r] & \pi_0\!\big(R[T]^{\otimes_R p}\big)^{tC_p} \arrow[r, equal] & \pi_0\!\big(R[T^p]\big)^{tC_p} \\
& & & \pi_0\!\big(R[T]\big)^{tC_p} \arrow[u, "\Delta"', "\sim"] \arrow[r, equal] & R/p\,[T]
\end{tikzcd}
\end{equation}
Here the first $\otimes$ is over $\mathbb Z$ (the uncompleted, section-wise tensor
of sheaves of abelian groups), the middle arrow is the base change to
$\otimes_R$, and the K\"unneth identity for free modules gives
$R[T]^{\otimes_R p}=R[T^p]$ with $C_p$ acting by permuting the factors of $T^p$.
The diagonal embedding $T\hookrightarrow T^p$ induces a $C_p$-equivariant map
$\Delta\colon R[T]\to R[T^p]$ (trivial action on the source), hence
$\pi_0(R[T])^{tC_p}=R/p\,[T]\to\pi_0(R[T^p])^{tC_p}$.

\begin{claim}
\label{claim:13:delta-iso}
$\Delta$ induces an isomorphism $\pi_0(R[T])^{tC_p}\xrightarrow{\ \sim\ }\pi_0(R[T^p])^{tC_p}$.
\end{claim}

Granting \cref{claim:13:delta-iso}, the composite in \eqref{eq:13:frobenius-RT}
lands isomorphically in $\pi_0(R[T])^{tC_p}=R/p\,[T]$ (trivial action, so the Tate
construction is reduction mod $p$), producing the required Frobenius-semilinear
map $R[T]\to R/p\,[T]$, functorial in $T$ and compatible with maps to the ground
ring. This proves \cref{thm:13:frobenius}.

\begin{remark}[What is new in the light setting]
\label{rmk:13:new-in-light}
In the previous iteration of this material \cref{claim:13:delta-iso} had to be
imposed as an extra \emph{axiom} on analytic rings. In the light-profinite
setting it can instead be \emph{proved} for arbitrary $R$; that is the point of
the following lemma.
\end{remark}

\Cref{claim:13:delta-iso} is the case $S=T^p$ (with $S^{C_p}=T$, the diagonal) of:

\begin{lemma}[Fixed points and Tate cohomology]
\label{lem:13:tate-fixed-points}
Let $S$ be a light profinite set with a $C_p$-action, and $S^{C_p}\subseteq S$
its (closed, again light profinite) subset of fixed points. Then the inclusion
$R[S^{C_p}]\hookrightarrow R[S]$ induces an isomorphism on Tate cohomology
$\widehat H^{*}(C_p,-)$, in particular on $\pi_0$.
\end{lemma}

\begin{proof}
\emph{Step 1: the inclusion is injective, and it suffices to kill the quotient.}
Any inclusion of light profinite sets admits a retraction
(\cref{prop:2:injective}); applying it to $S^{C_p}\hookrightarrow S$ and hitting
with $R[-]$ shows $R[S^{C_p}]\hookrightarrow R[S]$ is (split) injective. By the
long exact sequence in Tate cohomology it therefore suffices to show that
\begin{equation}
\label{eq:13:quotient-vanishes}
R[S]\big/R[S^{C_p}]\quad\text{has vanishing Tate cohomology.}
\end{equation}

\emph{Step 2: the quotient depends only on the complement.} For $T\subseteq S$ a
closed light profinite subset, $R[S]/R[T]$ depends only on the (locally light
profinite, i.e.\ locally compact) complement $S\setminus T$. Precisely: given
$S'\to S$ with pullback $T'=S'\times_S T$, if $S'\to S$ is an isomorphism over
$S\setminus T$ then
\begin{equation}
\label{eq:13:quotient-independence}
R[S']\big/R[T']\ \xrightarrow{\ \sim\ }\ R[S]\big/R[T].
\end{equation}
This holds because the square
% board 34 @ 01:17:32 --- modifying (compactifying) T without changing the complement; a pushout of condensed sets
\begin{equation}
\label{eq:13:pushout-square}
\begin{tikzcd}
T' \arrow[r, hook] \arrow[d] & S' \arrow[d, "\text{iso over }S\setminus T"] \\
T \arrow[r, hook] & S
\end{tikzcd}
\end{equation}
is a pushout in (light) condensed sets --- a small exercise with the Grothendieck
topology. Intuitively one is free to replace $T$ by any larger closed set, i.e.\
to choose a different compactification of $S\setminus T$.

\emph{Step 3: free actions on such spaces are induced.} If $X$ is a
$\sigma$-compact totally disconnected locally compact Hausdorff space
($\sigma$-compact $=$ countable union of compacts) on which $C_p$ acts with no
fixed points, then
\begin{equation}
\label{eq:13:induced}
X\ \cong\ \coprod_{C_p} Y,\qquad C_p\ \text{permuting the copies.}
\end{equation}
Indeed $X\to X/C_p$ is a covering map (by local compactness, Hausdorffness and
freeness of the action), and $X/C_p$ is again $\sigma$-compact, totally
disconnected, locally compact Hausdorff, hence a countable disjoint union of
light profinite sets. Over each such profinite piece the cover is trivial (total
disconnectedness makes it locally, hence globally, trivial), and trivializing
gives \eqref{eq:13:induced} with $Y=X/C_p$.

\emph{Conclusion.} Combining Steps 2 and 3: the complement $S\setminus S^{C_p}$ is
locally light profinite with a free $C_p$-action (a light profinite set minus a
light closed subset is $\sigma$-compact totally disconnected LCH, by lightness),
so it is $\coprod_{C_p}Y$; recompactifying by compactifying $Y$ and taking the
$C_p$ copies, Step 2 identifies
\begin{equation}
\label{eq:13:quotient-induced}
R[S]\big/R[S^{C_p}]\ \cong\ \bigoplus_{C_p} X',
\end{equation}
an \emph{induced} $C_p$-module. Induced modules have vanishing Tate cohomology,
giving \eqref{eq:13:quotient-vanishes}.
\end{proof}

\begin{remark}[Frobenius, symmetric powers, and animated analytic rings]
\label{rmk:13:symmetric-powers}
\Cref{thm:13:frobenius} is not merely cute. Suppose $(\ur R,D_{\ge 0}(R))$ is a
derived pre-analytic ring and we equip $\ur R$ with the structure of an
\emph{animated} (condensed) commutative ring. This produces the derived symmetric
power functors $\Sym^i$ on $D_{\ge 0}(\ur R)$ --- the functors that build free
animated rings, i.e.\ the monad describing animated rings over $\ur R$. The
Frobenius theorem implies that the $\Sym^i$ descend to the localization
$D_{\ge 0}(R)$: if $M\to N$ becomes an isomorphism in $D_{\ge 0}(R)$ (i.e.\ after
completion), then so does $\Sym^i M\to\Sym^i N$, so $\Sym^i$ is defined on
$D_{\ge 0}(R)$. The deduction from the Frobenius statement uses Dold--Puppe's
computation of the stable derived functors of the symmetric power functors
\cite{Dold_1961}; details are in \cite{clausen2020analytic}. Consequently one can
carry out the completion of \cref{prop:13:completion} in the animated context as
well, yielding a good category of \emph{animated analytic rings}: pairs
$(\ur R,D_{\ge 0}(R))$ with $\ur R$ an animated condensed ring and $D_{\ge 0}(R)$
an analytic ring structure on the underlying $\mathbb{E}_\infty$-ring.
\end{remark}

\subsection{Killing an algebra object}

The last topic is a general mechanism, of which solidification is the motivating
case. Recall that solid $\mathbb Z$ was presented as the analytic ring structure
in which $M\in\Solid_{\mathbb Z}$ iff
$\iHom(P,M)\xrightarrow{\,1-\sigma\,}\iHom(P,M)$ is an isomorphism, where
$P=\mathbb Z[\mathbb N\cup\{\infty\}]/[\infty]$ and $\sigma$ is the shift. This is
equivalent to the vanishing
\begin{equation}
\label{eq:13:solid-as-vanishing}
\RHom(A,M)=0,\qquad A\ :=\ P/(T-1),
\end{equation}
where $T=\sigma$ acts via the ring structure on $P$ (measures on $\mathbb N$, with
multiplication from addition; a priori $A$ is the derived quotient, but
multiplication by $T-1$ is injective, and in any case the class of modules is
unchanged). Thus $A$ is an algebra object of the ambient category
$D(\CondAbl)$, and one is ``killing'' it.

\begin{definition}[General context]
\label{def:13:killing}
Let $(\mathcal C,\otimes)$ be a presentable stable symmetric monoidal
$\infty$-category with $\otimes$ commuting with colimits in each variable. Let
$A\in\mathcal C$ be an ``algebra object'' in a very weak sense: a multiplication
$\mu\colon A\otimes A\to A$ and a unit $\eta\colon \mathbf 1\to A$ such that the
composite
\begin{equation}
\label{eq:13:weak-unital}
A\ \xrightarrow{\ \eta\otimes A\ }\ A\otimes A\ \xrightarrow{\ \mu\ }\ A
\end{equation}
is an isomorphism (one-sided unitality). No associativity is required. Define
\begin{equation}
\label{eq:13:killed-subcat}
\mathcal D\ \subseteq\ \mathcal C,\qquad \mathcal D=\{M\in\mathcal C\ :\ \iHom(A,M)=0\},
\end{equation}
the objects for which $A$ is ``internally $0$''. The goal is, in favorable
situations, to give a formula for the left adjoint to the inclusion
$\mathcal D\subseteq\mathcal C$.
\end{definition}

\begin{example}[Three instances]
\label{ex:13:killing-examples}
\begin{enumerate}
\item $\mathcal C=D(\CondAbl)$, $A=P/(T-1)$: then $\mathcal D=D(\Solid_{\mathbb Z})$;
similarly $\Solid_{\mathbb Z[T]}$ arises by killing the corresponding algebra.
\item $\mathcal C=D(R)$ for an ordinary ring $R$, $A=R/f$ (the cofiber of
$\cdot f$): then $\mathcal D=D(R[1/f])$, i.e.\ inverting $f$.
\item $\mathcal C=D(R)$, $A=R[f^{-1}]=R[1/f]$: then $\mathcal D=D(R)^{\wedge}_f$,
the $f$-complete objects, and the left adjoint is derived $f$-completion.
\end{enumerate}
\end{example}

\begin{construction}[The functor $F$ and its iteration]
\label{constr:13:F}
Let $C=\mathrm{fib}(\mathbf 1\to A)$ be the fiber of the unit, and define
\begin{equation}
\label{eq:13:F-def}
F\colon\mathcal C\to\mathcal C,\qquad F(X)=\iHom(C,X).
\end{equation}
The map $C\to\mathbf 1$ induces a natural transformation
\begin{equation}
\label{eq:13:X-to-FX}
X=\iHom(\mathbf 1,X)\ \longrightarrow\ \iHom(C,X)=F(X).
\end{equation}
Iterating --- applying $F$ to the previous \emph{map} (not the instance of the
previous map at $F(X)$) --- gives a sequence
\begin{equation}
\label{eq:13:iterate}
X\ \longrightarrow\ F(X)\ \xrightarrow{\ F(-)\ }\ F^2(X)\ \longrightarrow\ F^3(X)\ \longrightarrow\ \cdots,
\qquad
F^\infty(X):=\varinjlim_n F^n(X).
\end{equation}
\end{construction}

\begin{claim}
\label{claim:13:hom-invariance}
If $M\in\mathcal D$, then $\iHom(-,M)$ applied to $X\to F(X)$ is an isomorphism.
\end{claim}

\begin{proof}
It suffices that the fiber of $X\to F(X)$ be an $A$-module. Since
$A=\mathrm{cofib}(C\to\mathbf 1)$, the fiber of
$\iHom(\mathbf 1,X)\to\iHom(C,X)$ is, up to a shift, $\iHom(A,X)$, which carries
an $A$-module structure (acting on the $A$-variable). Now any $A$-module is a
retract of $A\otimes(\text{itself})$, by the unitality axiom
\eqref{eq:13:weak-unital}; and $\iHom(A\otimes-,M)=\iHom(-,\iHom(A,M))=0$ because
$M\in\mathcal D$. Hence homing out of an $A$-module into $M$ vanishes, so
$\iHom(-,M)$ sees $X\to F(X)$ as an isomorphism.
\end{proof}

\begin{proposition}[The left adjoint]
\label{prop:13:left-adjoint}
Suppose either
\begin{enumerate}
\item the colimit \eqref{eq:13:iterate} stabilizes for all $X$ (every map from
some point on is an isomorphism), or
\item $\iHom(A,-)$ commutes with (sequential) colimits.
\end{enumerate}
Then $X\mapsto F^\infty(X)$ is the left adjoint to the inclusion
$\mathcal D\subseteq\mathcal C$.
\end{proposition}

\begin{proof}
Two things must be checked.

\emph{(a) $\Hom(X,-)\simeq\Hom(F^\infty X,-)$ on $\mathcal D$.} Just as for
$X\to F(X)$, each map $F^n X\to F^{n+1}X$ is an isomorphism on $\iHom(-,M)$ for
$M\in\mathcal D$: applying $F$ preserves the property used in
\cref{claim:13:hom-invariance} (the fiber is an $A$-module, and internally homing
into an $A$-module gives an $A$-module). Homing out of the colimit is the inverse
limit of homing out of the terms, so $X\to F^\infty X$ is an isomorphism on
$\iHom(-,M)$; here one uses hypothesis (1) or (2) to pull the internal Hom inside
the colimit.

\emph{(b) $F^\infty X\in\mathcal D$.} One checks $\iHom(A,F^\infty X)=0$. Under
the hypothesis $\iHom(A,-)$ commutes with the colimit, and the transition maps
$\iHom(A,F^n X)\to\iHom(A,F^{n+1}X)$ are null-homotopic (internally homing $A$
shifts one step along the sequence), so the colimit vanishes.
\end{proof}

\begin{remark}[Recovering the classical formulas]
\label{rmk:13:classical-formulas}
The formula $F^\infty$ specializes to familiar left adjoints
(\cref{ex:13:killing-examples}).
\begin{itemize}
\item For $A=R/f$ (inverting $f$) the left adjoint is
\begin{equation}
\label{eq:13:invert-f}
L(M)\ =\ \varinjlim\big(M\xrightarrow{f}M\xrightarrow{f}M\xrightarrow{f}\cdots\big)\ =\ M[1/f],
\end{equation}
which $F^\infty$ reproduces as a sequential colimit.
\item For $A=R[1/f]$ ($f$-completion) the left adjoint is
\begin{equation}
\label{eq:13:f-completion}
L(M)\ =\ \varprojlim_n M/f^n\ =\ F(M),
\end{equation}
where the iteration stabilizes after the first term: the inverse limit is
already seen at $F(M)$, and the subsequent maps are isomorphisms.
\item For solidification, $F(X)$ is explicitly a quotient of null sequences in
$X$ modulo the telescoping relations (those forcing a null sequence to be summed);
iterating computes the solidification of a general module, and one can check it
recovers the classical formulas.
\end{itemize}
\end{remark}

\begin{remark}[When the colimit stabilizes]
\label{rmk:13:idempotent}
The stabilization hypothesis (1) of \cref{prop:13:left-adjoint} holds, for
example, when $A$ is idempotent --- as in \cref{ex:13:killing-examples}(3),
$A=R[1/f]$. By contrast $A=R/f$ in \cref{ex:13:killing-examples}(2) is
\emph{not} idempotent, and there one uses hypothesis (2), that $\iHom(A,-)$
commutes with colimits. (Solidifying $R[1/f]$ would instead be the other
operation, not the algebraic inverting of $f$.)
\end{remark}
